\documentclass{article}
\usepackage[utf8]{inputenc}
\usepackage[russian]{babel}
\usepackage{amsmath}
\usepackage{amsfonts}
\usepackage{array}

\begin{document}

\title{Домашнее задание №3 по обработке сигналов}
\author{}
\date{}
\maketitle

\section{Задача 1}

Дана порождающая матрица
\[
G = \begin{pmatrix}
1 & 1 & 0 & 1 & 0 \\
0 & 1 & 1 & 0 & 1 \\
1 & 0 & 1 & 1 & 0
\end{pmatrix}.
\]

Код линейный двоичный с параметрами $n=5$, $k=3$, $d_{\min}=1$ (см. вычисление ниже).

\textbf{Проверочная матрица $H$} (одна из возможных, строки образуют базис ядра $G$ над $\mathbb{F}_2$):
\[
H = \begin{pmatrix}
0 & 1 & 1 & 1 & 0 \\
1 & 0 & 0 & 1 & 0
\end{pmatrix}.
\]

Проверка: $G H^T = 0$ (по модулю 2).

\textbf{Характеристики кода:}
\begin{itemize}
\item Длина: $n=5$
\item Размерность: $k=3$
\item Число кодовых слов: $2^3=8$
\item Минимальное расстояние: $d_{\min}=1$
\item Способность исправления ошибок: $t = \lfloor (d_{\min}-1)/2 \rfloor = 0$
\item Скорость кода: $R = k/n = 3/5$
\end{itemize}

Для нахождения $d_{\min}$ и всех кодовых слов использовался следующий Python-код:

\begin{verbatim}
import numpy as np
from itertools import product

G = np.array([[1,1,0,1,0],
              [0,1,1,0,1],
              [1,0,1,1,0]], dtype=int) % 2

weights = []
for coeffs in product([0,1], repeat=3):
    u = np.array(coeffs, dtype=int)
    c = (u @ G) % 2
    wt = np.sum(c)
    weights.append(wt)

print("Минимальное расстояние:", min(w for w in weights if w > 0))
\end{verbatim}

Результат: $d_{\min}=1$ (существует кодовое слово веса 1: $u=(1,1,1) \to c=(0,0,0,0,1)$).

\textbf{Таблица синдромного декодирования}

Синдром $s = H r^T \pmod{2}$ (вектор-столбец длины 2).  
Для каждого синдрома выбран лидер косета минимального веса (вычислено перебором всех $2^5=32$ векторов).

\begin{center}
\begin{tabular}{|c|c|c|}
\hline
Синдром $s$ & Лидер ошибки $e$ & Вес \\
\hline
$(0,0)$ & $(0,0,0,0,0)$ & 0 \\
$(0,1)$ & $(1,0,0,0,0)$ & 1 \\
$(1,0)$ & $(0,0,1,0,0)$ & 1 \\
$(1,1)$ & $(0,0,0,1,0)$ & 1 \\
\hline
\end{tabular}
\end{center}

(Примечание: ошибки в позициях 2 и 5 не выбраны лидерами, так как имеют те же синдромы, что и уже выбранные векторы веса 1 или 0.)

Для построения таблицы использовался аналогичный Python-код с перебором и выбором минимального веса для каждого синдрома.

\section{Задача 2}

Двоичный код Голея — совершенный линейный код с параметрами $[23,12,7]$.

Минимальное расстояние $d_{\min}=7$.

\textbf{Доказательство достижения границы Хэмминга (сферной упаковки):}

Код исправляет $t = \lfloor (7-1)/2 \rfloor = 3$ ошибки.

Размер сферы радиуса 3 в пространстве $\mathbb{F}_2^{23}$:
\[
V = \sum_{i=0}^{3} \binom{23}{i} = \binom{23}{0} + \binom{23}{1} + \binom{23}{2} + \binom{23}{3}.
\]

Вычисление:
\[
\binom{23}{0}=1, \quad \binom{23}{1}=23, \quad \binom{23}{2}=\frac{23\cdot22}{2}=253, \quad \binom{23}{3}=\frac{23\cdot22\cdot21}{6}=1771.
\]
\[
V = 1 + 23 + 253 + 1771 = 2048 = 2^{11}.
\]

Число кодовых слов $M = 2^{12}$.

Объём всех сфер:
\[
M \cdot V = 2^{12} \cdot 2^{11} = 2^{23},
\]
что в точности равно размеру всего пространства $\mathbb{F}_2^{23}$.

Следовательно, сферы радиуса 3 вокруг кодовых слов заполняют всё пространство без пересечений — код совершенный и достигает границы Хэмминга с равенством.

(Для подтверждения суммы использовался Python-код ниже.)

\begin{verbatim}
from scipy.special import binom
V = binom(23,0) + binom(23,1) + binom(23,2) + binom(23,3)
print(int(V))  # 2048
\end{verbatim}

\end{document}